\documentclass[12pt]{article}
\usepackage[utf8]{inputenc}
\usepackage[T1]{fontenc}
\usepackage[margin=1in]{geometry}
\usepackage{setspace}
\usepackage{xcolor}
\usepackage[colorlinks=true, urlcolor=blue!60!black, linkcolor=blue!60!black]{hyperref}
\setlength{\parindent}{2em}
\setlength{\parskip}{1.25em}
\renewcommand{\baselinestretch}{1.0}
\begin{document}
\begin{center}
{\LARGE \textit{Frontiero v. Richardson}, 411 U.S. 677 (1973)\par}
\vspace{0.8em}
{\large Case Brief --- Four-Page Edition\par}
\end{center}
\vspace{1em}

\section{Memory Jogger}

In an 8--1 judgment, the Court invalidated military benefit statutes that automatically treated servicemen's wives as dependents while requiring servicewomen to prove their husbands' actual dependency; a four-Justice plurality famously declared sex a suspect classification subject to strict scrutiny (\href{https://supreme.justia.com/cases/federal/us/411/677/}{Justia}).

\section{Detailed Case Facts}

Air Force Lieutenant Sharron Frontiero sought increased housing allowances and medical and dental benefits for her husband, Joseph, as a dependent. Under 37 U.S.C. \S\S~401, 403 and 10 U.S.C. \S\S~1072, 1076, a male service member could claim his wife as a dependent automatically, but a female member had to prove that her husband depended on her for more than one-half of his support. Joseph, a full-time college student, had about \$354 in monthly expenses and received \$205 a month in veterans' benefits, so his wife could not make that showing and the Air Force denied her application. A male officer with an identically situated wife would have received the benefits automatically.

\section{Procedural History}

The Frontieros sued the Secretary of Defense in the Middle District of Alabama. A three-judge court upheld the statutes 2--1 under rational-basis review: because husbands were typically breadwinners and about 99\% of service members were male, Congress could presume wives' dependency to save administrative expense. Judge Johnson dissented under \textit{Reed v. Reed}. On direct appeal, Ruth Bader Ginsburg made her first Supreme Court argument, as amicus for the ACLU.

\section{Judicial Votes}

Reversed, \textbf{8--1}, with no majority opinion on the standard of review (decided May 14, 1973).
\begin{itemize}
\item \textbf{Plurality:} Justice Brennan, joined by Justices Douglas, White, and Marshall.
\item \textbf{Concurring in the judgment:} Justice Stewart; Justice Powell, joined by Chief Justice Burger and Justice Blackmun.
\item \textbf{Dissent:} Justice Rehnquist.
\end{itemize}

\section{Holding}

The statutes violate the Fifth Amendment's Due Process Clause insofar as they require servicewomen to prove a spouse's dependency to obtain benefits servicemen receive automatically. The binding rule is the narrowest ground: under \textit{Bolling v. Sharpe} and \textit{Reed}, statutes that treat similarly situated men and women differently solely for administrative convenience are unconstitutional. The remedy extended automatic benefits to servicewomen. Brennan's declaration that sex is inherently suspect drew only four votes; \textit{Craig v. Boren} (1976) later settled on intermediate scrutiny.

\section{Analysis of Opinions}

\textbf{Plurality (Brennan).} The statutes imposed a procedural burden of proof only on servicewomen and substantively denied them benefits that servicemen receiving the same support obtained. Reading \textit{Reed} as a departure from rational-basis review, Brennan treated sex like race, alienage, and national origin, citing a long history of romantic paternalism (as in Justice Bradley's \textit{Bradwell} concurrence), the visibility of sex and women's underrepresentation in political institutions, sex as an immutable accident of birth, its frequent irrelevance to ability, and Congress's own recent turn against sex discrimination in the Equal Pay Act, Title VII, and the Equal Rights Amendment. Under strict scrutiny, administrative convenience failed: the government offered no proof that automatic benefits were cheaper than affidavits, and ``the Constitution recognizes higher values than speed and efficiency.''

\textbf{Stewart (concurring in the judgment).} The statutes ``work an invidious discrimination in violation of the Constitution'' under \textit{Reed}.

\textbf{Powell (concurring in the judgment).} \textit{Reed} sufficed. With the ERA before the states, declaring sex a suspect class would preempt a major political decision and weaken confidence in judicial restraint.

\textbf{Rehnquist (dissenting).} In one sentence, he adopted Judge Rives's opinion upholding the scheme as a rational cost-saving generalization.

\section{Examples: Future Applications}

\begin{enumerate}
\item \textbf{Same-Side (Invalid): Survivors' Benefits.} Social Security pays widows automatically but requires widowers to prove half-support from their wives. Like \textit{Frontiero}, it rests on breadwinner stereotypes (compare \textit{Califano v. Goldfarb}).
\item \textbf{Same-Side (Invalid): Parental Leave.} An agency grants paid caregiver leave automatically to mothers but makes fathers prove they are the primary caregiver.
\item \textbf{Opposite-Side (Valid): Uniform Dependency Test.} Congress requires \emph{all} service members to show by affidavit that their spouse earns less than half the household's expenses. The test turns on need, not sex.
\item \textbf{Opposite-Side (Valid): Disability Documentation.} A state requires medical documentation of disability for rehabilitation subsidies, waiving it for those already certified. Disability, unlike sex, bears on functional needs.
\item \textbf{Fence-Sitter: Widows' Tax Exemption.} A state grants a property-tax exemption to widows but not widowers, citing women's wage disparities. Footnote 22 distinguished laws that rectify past discrimination; strict scrutiny would still be skeptical, while a deferential remedial approach (\textit{Kahn v. Shevin}) might uphold it.
\end{enumerate}

\section{Critique}

\textbf{Originalist critiques.} \textit{Frontiero} rests on reverse incorporation of equal protection into a Fifth Amendment that has no such clause, and neither the Fifth nor the Fourteenth Amendment was originally understood to make sex suspect; Brennan treated history as a record to be overcome. On this view Powell's deference to the Article V process was institutionally sound.

\textbf{Progressive and feminist critiques.} Scholars praise the dismantling of romantic paternalism but lament that the Court fell one vote short of strict scrutiny because of Powell's wait-and-see posture on an ERA that never was ratified, leaving intermediate scrutiny after \textit{Craig}. Reva Siegel adds that the compromise obscured the Nineteenth Amendment's transformation of women's citizenship, and disability scholars fault the dicta calling disability nonsuspect.

\section{Key Quotations}

\begin{quote}
``And what differentiates sex from such nonsuspect statuses as intelligence or physical disability, and aligns it with the recognized suspect criteria, is that the sex characteristic frequently bears no relation to ability to perform or contribute to society.''
\end{quote}
\noindent --- Justice Brennan, plurality opinion, 411 U.S. at 686--87.

\begin{quote}
``There can be no doubt that our Nation has had a long and unfortunate history of sex discrimination. Traditionally, such discrimination was rationalized by an attitude of `romantic paternalism' which, in practical effect, put women not on a pedestal, but in a cage.''
\end{quote}
\noindent --- Justice Brennan, plurality opinion, 411 U.S. at 684.

\begin{quote}
``With these considerations in mind, we can only conclude that classifications based upon sex, like classifications based upon race, alienage, or national origin, are inherently suspect, and must therefore be subjected to strict judicial scrutiny.''
\end{quote}
\noindent --- Justice Brennan, plurality opinion, 411 U.S. at 688.

\end{document}
